% ============================================================================
%  Chapitre 5, partie I : Position du problème
% ============================================================================

% ----------------------------------------------------------------------------
\section{Position du problème}
% ----------------------------------------------------------------------------

La simulation suppose connus la masse $m$, le rayon $a$ et la matrice
d'inertie $\mat{I}_b$ de la boule. Pour une vraie boule, $m$ et $a$ se
mesurent directement, avec une balance et un pied à coulisse. Ce n'est
pas le cas de $\mat{I}_b$, qui dépend de la répartition de la masse à
l'intérieur, ce que l'on ne connaît pas lors de la production. Le but de ce chapitre est de répondre à la question : "Comment peut-on les déterminer à partir de mesures embarquées ?". Dans ce chapitre, on utilise un gyromètre au centre de la boule, et éventuellement un accéléromètre.

\subsection{Le repère du gyromètre}

Dans les chapitres précédents, on a écrit l'inertie dans les axes
principaux de la boule, où $\mat{I}_b = \mathrm{diag}(I_1, I_2, I_3)$.
\textbf{Cependat, on ne connaît jamais ces axes}, et on ne peut donc pas y aligner le gyromètre. On doit travailler alors dans le repère du gyromètre, relié aux axes principaux par une rotation $\mat{P}$ inconnue. Dans ce repère, la matrice d'inertie
\[
    \mat{J} = \mat{P}\,\mat{I}_b\,\mat{P}^T
    = \begin{pmatrix}
        J_{11} & J_{12} & J_{13}\\
        J_{12} & J_{22} & J_{23}\\
        J_{13} & J_{23} & J_{33}
    \end{pmatrix}
\]
qui n'est plus diagonale, mais reste symétrique définie positive par définition. Estimer l'inertie revient donc à estimer ses $6$ coefficients, on ne peut pas faire autrement sans connaître $\mat{P}$.

\subsection{Retour aux axes principaux}

Une fois $\mat{J}$ estimée, on la diagonalise. Elle est symétrique,
donc
\[
    \mat{J} = \mat{P}\,\mathrm{diag}(I_1, I_2, I_3)\,\mat{P}^T,
    \qquad \mat{P} \in \SO
\]
Les valeurs propres de $\mat{J}$ sont les moments principaux
$I_1, I_2, I_3$, et les colonnes de $\mat{P}$ donnent les axes
principaux exprimés dans le repère du gyromètre. Donc, en plus, en diagonalisant $\mat{P}$, on obtient aussi la position du gyromètre par rapport aux axes principaux de l'inertie.


